\documentclass[10pt,a4paper]{amsart}
\usepackage[margin=19mm]{geometry}
\usepackage{amsmath,amssymb,mathtools}
\usepackage[scr=rsfs,cal=euler]{mathalfa}
\usepackage{xfrac}
\usepackage[unicode,hidelinks,bookmarks=false]{hyperref}
\newcommand{\ie}{i.e.\ }
\newcommand{\cf}{cf.\ }
\newcommand{\resp}{respectively\ }
\newcommand{\Subsection}{\subsection}
\providecommand{\nobreakdash}{\nobreak\hbox{-}}
\setlength{\emergencystretch}{2em}
\multlinegap0pt
\usepackage{amssymb}
\usepackage{mathtools}
\usepackage{tikz}
\usepackage{float}
\usepackage{xcolor}
\newtheorem{theorem}{Theorem}
\newtheorem{corollary}{Corollary}
\newtheorem{lemma}{Lemma}
\newcommand{\B}{B_1}
\newcommand{\RR}{\mathbb R}
\newcommand{\Rot}{\mathcal R}
\newcommand{\Sph}{\mathbb{S}^2}
\newcommand{\eiii}{e_3}
\title{Br\'ezis-Coron Uniqueness of Harmonic Maps via Boundary Rigidity of Hopf Differentials}
\author{Qi Guo, Xueping Huang, Yi C. Huang, Fanghua Lin,, Juncheng Wei}
\date{Source: arXiv:2606.08648v1 (CC BY 4.0)}
\begin{document}
\maketitle

\noindent\emph{Source and licence.} Br\'ezis-Coron Uniqueness of Harmonic Maps via Boundary Rigidity of Hopf Differentials. Version arXiv:2606.08648v1 (CC BY 4.0), distributed by arXiv under the Creative Commons Attribution 4.0 International licence, http://creativecommons.org/licenses/by/4.0/, as recorded in the arXiv licence metadata for this exact version and shown on the abstract page https://arxiv.org/abs/2606.08648v1. Full licence notice: \texttt{licenses/arxiv-2606.08648-CC-BY-4.0.txt}.\par\smallskip

\noindent\textit{Editorial scope.} Counted unit: the article's corollary cor:normal-trace-Hsystem (source0.tex lines 289--302). The article's complete original proof of it is delivered with it (source0.tex lines 1433--1556, one \texttt{proof} environment of 124 lines, which the article places away from the statement). No mathematics is written, repaired or re-derived: the statement and the proof are the article's own lines, in the article's own notation, and no retained line is substituted or re-flowed.  Retained uncounted as the source-local context the proof needs: lines 233--235; lines 1318--1330.  Nothing was omitted from any retained span, so the package carries no omission record.  No retained line contains a citation, so the wrapper carries no bibliography and no bibliography entry.  No retained line contains a figure environment, an image-inclusion command or drawing code, so no image is delivered and the package declares no asset dependency.  Editorial formatting only, in addition: an AMS \texttt{amsart} preamble that restates the article's own declarations and macros used by the retained lines, namely the environments \texttt{theorem}, \texttt{corollary}, \texttt{lemma} on separate counters, together with the standard packages those lines need.  The retained environments print in this document as Theorem 1, Corollary 1, Lemma 1 in order of appearance, whereas the article prints the same items with the numbering of its own full document.  The complete native source of the article as distributed in the arXiv e-print for this exact version is bundled unchanged as \texttt{sources/202346/source0.tex}.
\begin{theorem}\label{thm:main}
Let $0<R<1$ and let $u\in H^1(B_1;\Sph)$ be a weak harmonic map with trace $g_R$. Then $u=u^+$ or $u=u^-$.
\end{theorem}

\begin{lemma}\label{reverseimpossible}
Let $Y:B_1\to\mathbb R^3$ be a smooth conformal immersion satisfying the
$H$-system
\[
 -\Delta Y=Y_x\wedge  Y_y \quad\text{in }B_1.
\]
Let $\nu_Y$
be the unit normal induced by the orientation of the parameter disk.  Then  $\nu_Y$ cannot be
 the orientation-reversing alternative, i.e.
\[
 \nu_Y=u^-\circ\mathcal R_\phi.
\]
\end{lemma}

\begin{corollary}\label{cor:normal-trace-Hsystem}
Let $Y\in C^2(\B;\RR^3)\cap C^1(\overline{\B};\RR^3)$ be a conformal immersion satisfying
\begin{align*}\displaystyle
\begin{cases}
        -\Delta Y=Y_x\wedge Y_y,\quad\text{in }\B.\\
        \nu_Y(1,\theta)=g_R^\sigma(\theta+\phi),
\end{cases}
\end{align*}
for some $\sigma\in\{+1,-1\}$ and some $\phi\in\RR$.
Then there exists $a\in\RR^3$ such that
\[
Y=a+Z^\sigma\circ\Rot_\phi.
\]
\end{corollary}

\begin{proof}[Proof of Corollary~\ref{cor:normal-trace-Hsystem}]
We give the proof in three steps.

\smallskip
\noindent\textbf{Step 1: the induced normal is a harmonic map.}

Write the conformal factor as
\[
|Y_x|^2=|Y_y|^2=e^{2\lambda},
\qquad
Y_x\cdot Y_y=0.
\]
Since $Y$ is oriented by $(x,y)$,
\[
Y_x\wedge Y_y=e^{2\lambda}\nu_Y.
\]
For a conformal immersion into $\RR^3$, the mean-curvature scalar $H$ with respect to $\nu_Y$ satisfies
\[
\Delta Y=2e^{2\lambda}H\nu_Y.
\]
As in the proof of Lemma \ref{reverseimpossible}, the $H$-system equation gives $H=-1/2$.
In particular, the mean curvature is constant. The standard Gauss-map equation for a conformal immersion is
\[
\Delta_g\nu_Y+|A|_g^2\nu_Y=\nabla_g H,
\]
where $g=e^{2\lambda}(dx^2+dy^2)$ is the induced metric and $A$ is the second fundamental form. Since $H$ is constant, the right-hand side
vanishes. Multiplying by $e^{2\lambda}$ gives
\[
\Delta \nu_Y+|\nabla \nu_Y|^2\nu_Y=0.
\]
Thus $\nu_Y$ is a weak harmonic map from $\B$ into $\Sph$.

\smallskip
\noindent\textbf{Step 2: apply the Br\'ezis-Coron uniqueness theorem to the normal.}

(1) First, we suppose $\sigma=+1$. Define
\[
\widetilde\nu(r,\theta)=\nu_Y(r,\theta-\phi).
\]
Then $\widetilde\nu$ is a harmonic map into $\Sph$ and
\[
\widetilde\nu(1,\theta)=g_R^+(\theta).
\]
By Theorem~\ref{thm:main}, we have
\[
\widetilde\nu=u^+\quad\text{or}\quad \widetilde\nu=u^-.
\]
A direct computation gives
\[
u^+\cdot(u^+_x\times u^+_y)
=
\frac{4\rho^2}{(1+\rho^2r^2)^2}>0,
\qquad
u^-\cdot(u^-_x\times u^-_y)
=
-\frac{4\rho^2}{(\rho^2+r^2)^2}<0.
\]
Thus $u^+$ is conformal and orientation-preserving, whereas $u^-$ is conformal and orientation-reversing.

Lemma \ref{reverseimpossible} shows that the orientation-reversing alternative is incompatible with the $H$-system orientation. This excludes
$u^-$ in the case $\sigma=+1$, and gives
\[
\nu_Y=u^+\circ\Rot_\phi.
\]
Similar to Lemma \ref{reverseimpossible}, we have $H=-1/2$, $A=-I/2$. Then from the Weingarten formula $d\nu =-dY\circ A$, we have
\[ d\nu =\frac12 dY, \quad d Y=2 d \nu. \]
Since $B_1$ is connected, we have
\[
Y=a+2u^+\circ\Rot_\phi.
\]
Since $Z^+=2u^+-2\sqrt{1-R^2}\,\eiii$, the constant $2\sqrt{1-R^2}\,\eiii$ can be absorbed into $a$, and therefore
\[
Y=a+Z^+\circ\Rot_\phi.
\]

\smallskip
(2) Now suppose $\sigma=-1$. Set
\[
\widetilde\nu(r,\theta)=-\nu_Y(r,\theta-\phi+\pi).
\]
Then $\widetilde\nu$ is again a harmonic map into $\Sph$, and
\[
\widetilde\nu(1,\theta)
=-g_R^-(\theta+\pi)
=g_R^+(\theta).
\]
By Theorem~\ref{thm:main}, we have
\[
\widetilde\nu=u^+\quad\text{or}\quad \widetilde\nu=u^-.
\]
The first alternative would make $\nu_Y$ orientation-reversing, because the antipodal map on $\Sph$ reverses orientation. This is impossible
by the shape-operator argument as in Lemma \ref{reverseimpossible}. Therefore
\[
\widetilde\nu=u^-.
\]
Equivalently,
\[
\nu_Y=-u^-\circ\Rot_{\phi+\pi}.
\]
Hence, using again $Y=2\nu_Y+a$,
\[
Y=a-2u^-\circ\Rot_{\phi+\pi}.
\]
By the definition of $Z^-$,
\[
Z^-= -2u^-\circ\Rot_\pi+2\sqrt{1-R^2}\,\eiii,
\]
and the vertical translation is absorbed into $a$. Thus
\[
Y=a+Z^-\circ\Rot_\phi.
\]

% \smallskip
% \noindent\textbf{Step 3: normalization.}
% The maps $Z^\pm$ satisfy
% \[
% Z^\pm(1,\theta)=\Gamma_R(\theta).
% \]
% The $H$-system equation is invariant under translations in $\RR^3$ and rotations of the disk parameter. Therefore imposing
% \[
% Y(1,\theta)=\Gamma_R(\theta)
% \]
% fixes the translation and the boundary parameter rotation, leaving only the two normalized representatives $Z^+$ and $Z^-$.
\end{proof}

\end{document}
